\section{引言}\label{sec:intro}
交通预测为路网运行管理、出行规划与拥堵预警提供前瞻性信息。然而，预测系统能够获得的历史数据并不总是连续且完整：传感器逐步部署使部分位置缺少训练期记录，设备失效和通信异常又会改变运行期间的数据可用性。已有关于未观测位置预测和动态传感器网络的研究已认识到这些问题~\citep{mei2023uncertainty,zhou2025mogernn,zhao2026structure}。对于同一个预测位置，新近加入、持续观测和近期中断所提供的信息条件并不相同。本文将历史窗口内实际可用观测的位置、数量及其时间分布统称为\textbf{观测支持}。观测支持的变化不仅产生需要处理的空缺，也决定模型能从本地历史与关联位置获得怎样的预测依据。

围绕不完整交通观测，已有研究形成了互补的解决途径。缺失感知方法通过图循环重建、条件扩散或联合重建与预测，利用时空关联补充不可用信息~\citep{cini2022filling,liu2023pristi,zuo2023graph,yu2024ginar}；归纳与部分感知预测将预测范围扩展至缺少训练信号或部署期持续观测的位置~\citep{mei2023uncertainty,zhou2025mogernn,liu2025spatiotemporal,zhao2026structure}；自适应图学习及空间分块表示则改善空间关系和多尺度上下文的组织~\citep{yu2018spatio,wu2019graph,bai2020adaptive,jiang2023pdformer,fang2025efficient}。这些研究已经证明，缺失处理、空间信息传递与未观测位置预测能够协同工作。因而，进一步的关键并非判断模型是否能够接受缺失输入，而是明确：当历史支持在预测窗口内变化时，模型应如何相应调整信息利用方式。

这一问题在归纳预测中尤其突出。目标节点没有参与动态信号训练，并不意味着它在推理时始终没有观测；同样，训练期已有记录也不保证部署时始终拥有完整历史。现有方法分别强调训练信号覆盖、缺失恢复、传感器配置适应或空间交互，但这些能力并不直接回答“同一目标对应不同历史支持时，预测计算应如何变化”。即使可用观测数相同，近期集中、早期集中和分散分布的历史也可能对未来状态提供不同依据。本文因此研究\textbf{动态观测支持下的归纳时空预测}：在路网拓扑已知、部分节点动态信号不参与训练的条件下，利用变化的历史观测支持组织空间推理与预测上下文。该设置连接节点归纳与窗口内支持变化，强调固定模型参数下的预测行为。

我们的核心认识是，\textbf{观测支持应同时作为输入有效性的描述和预测计算的条件}。观测掩码标明哪些读数可用；在此基础上，观测角色与状态年龄可以为模型组织历史信息提供显式的结构化线索。这里的角色表示活跃（AAS）、新接入（NAS）、未观测（VS）或失效（FS）等状态，状态年龄描述接入或中断过程的进展。它们的作用在于提供一种可检验的归纳偏置，而不要求包含完整掩码之外的额外信息。沿这一思路，历史补全负责提供可用表示，空间传播决定如何交换关联信息，跨节点时间上下文则为预测过程提供条件。三者的协作，使模型能够在本地证据与空间证据之间建立依赖当前支持条件的信息利用方式。

据此，我们提出观测支持感知预测框架 ALFNet。\textbf{观测支持感知历史补全（SAHC）}借鉴图循环插补思想~\citep{cini2022filling}，利用可见观测及图关联构建两个分支共享的历史输入。\textbf{观测条件化图预测（ALFGraph）}借鉴图专家混合与扩散卷积循环网络（DCRNN）的循环结构~\citep{zhou2025mogernn,li2018diffusion}，由观测状态决定传播算子的组合。\textbf{空间上下文时间条件分支（ALFTemporal）}借鉴空间分区交互~\citep{fang2025efficient} 与两阶段注意力~\citep{zhang2023crossformer}，从跨节点历史生成条件表示，并通过特征级线性调制（FiLM）~\citep{perez2018film}作用于图循环预测过程。预测输出再经过一个由末步角色控制幅度的图修正。图预测与时间分支围绕共同的历史支持协作，最终形成单一预测输出。本文的重点是这一条件化组织方式，而非重新提出图循环、分块注意力或线性调制等基础算子。这一组织方式落实为三项可检验的设计选择：图传播的组合只由观测状态决定，跨节点上下文只在编码末状态注入一次，状态年龄显式描述接入与中断的进展（\ref{sec:3_2}节）；第\ref{sec:4_4}节分别以消融检验。

本文的贡献概括为三点：（1）围绕拓扑已知、训练信号覆盖不完整的路网，将窗口内变化的观测支持纳入归纳预测设置，并以零历史（Unseen）、持续观测（Steady）、逐步加入（Join）和历史中断（Fault）条件刻画其预测要求；（2）提出 ALFNet，将观测支持感知历史补全（SAHC）、观测条件化图预测（ALFGraph）与空间上下文时间条件分支（ALFTemporal）组织为协同预测框架；（3）在 METR-LA、PEMS-BAY、PEMS04 与 PEMS08 四个数据集上开展统一条件评价，并结合组件消融、配对观测干预和样本层面分析，考察总体收益、机制归因及支持结构对预测行为的影响。总体结果显示，与所列八个比较方法（共同协议下的适配版本）相比，ALFNet 在四个数据集的跨条件平均误差（LifeAvg）三项指标上均为最优；在逐步加入（Join）与历史中断（Fault）两类主条件下，这一优势在四个数据集上保持一致；组件消融表明，跨节点时间上下文与时间条件分支的贡献最大，观测状态条件化传播与状态年龄带来较小但方向基本一致的改进，支持条件本身的独立作用仍有待进一步隔离；配对干预与布局分析表明，这种收益与模型对支持条件变化的响应相关，这一响应是在预测目标与非干预输入保持不变时观察到的。收益主要出现在本地证据部分存在的条件；在本地历史完全缺失的零历史（Unseen）条件下，优势并不在所有数据集上成立（两个流量数据集上未成立），我们在实验中对该边界予以说明。
